\documentclass{report}
\usepackage{amsfonts}
\newcommand\R{{\mathbb R}}
\newcommand\Q{{\mathbb Q}}
\newcommand\N{{\mathbb N}}
\pagestyle{empty}
\begin{document}
\large
\centerline{\Huge Analisi Matematica}
\renewcommand{\thefootnote}{ } 
\footnote{\sl Avete 3 ore di tempo. Ogni esercizio
vale sei punti. La sufficienza si ottiene con un punteggio $\geq$
18. Solo le risposte chiaramente giustificate saranno prese in
considerazione.}
\renewcommand{\thefootnote}{\arabic{footnote}} 
\vskip.1cm
\centerline{\Large Terza Prova di Autovalutazione}
\vskip1cm
\begin{enumerate}
\item Si determini l'insieme degli $x\in[0,\pi]$ per cui
\[
1-\cos x\leq \sin x.
\]

\item Si calcoli, per ogni $n\in\N$,
\[
\sum_{i=1}^n i^2.
\]

\item Data la successione $\{a_n\}_{n\in\N}$ definita da: $a_1=3$,
$a_{n+1}=a_n/2$, si calcoli
\[
\sum_{i=1}^\infty a_i.
\]
\item Si studi la convergenza del limite
\[
\lim_{n\to\infty}n^2\left\{\sqrt{1+\frac 3{n^2}}-\sqrt{1-\frac
1{n^2}}\right\}.
\]
\item Si studi la convergena della serie
\[
\sum_{n=1}^\infty \left(1-\cos \frac 1n\right).
\]

\item Si studi la convergenza della serie
\[
\sum_{n=1}^\infty\frac {n!}{n^n}.
\] 


\end{enumerate}
\end{document}
